\section{The chart Hilbert function and the length of a mixed weight space}
\label{sec:deformation-chart}

We prove three assertions used in Section~\ref{sec:single-character}, namely
the choice of a nonzero trace representative on the chart, the
Hilbert function of the leading chart, and the local multiplicity
bound for the original equations. These assertions concern graded
algebras once the deformation presentation and its tangent
conditions have been constructed.

\subsection{Chart coordinates and leading monomials}

Keep the DVR \(\mathcal O\), its fraction field \(F\), and the
augmentation \(x_A:A\to\mathcal O\). The algebra \(A\) is
\(\mathbb Z\)-graded, finitely generated over \(A_0\) by homogeneous
elements, and \(x_A\) is zero on every nonzero weight. The cut chart
is an algebra homomorphism
\[
           \phi:A\longrightarrow F[[X_i:i\in J]]
\]
with \(J\) finite and weights \(w_i\in\{-2,-1,1,2\}\). Its constant
coefficient is \(x_A\), it carries homogeneous elements to series of
the same weight, and its linear terms separate the axes. The last
condition means that for every \(i\) some element maps to a series
whose coefficient of \(X_i\) is nonzero and whose other linear
coefficients vanish. Taking its homogeneous component of weight
\(w_i\) preserves these properties. There are \(r\ge1\) negative
unit-weight axes and \(r'+1\ge1\) positive unit-weight axes.

If \(\phi(a)=0\) for the weight-zero trace lift \(a\), it can be
changed by a weight-zero relation. Indeed, the coefficient of a chosen
weight-\(-1\) variable gives a linear functional on \(A_{-1}\)
which is nonzero by axis separation. It kills
\((\ker x_A)^{(2)}_{\rm sat}\) because products of two series with
zero constant term have zero linear coefficients, and coefficients
are taken in a field of fractions. The weight-\(-1\) equations
span the saturated cotangent lattice over \(F\). Consequently
\(\phi(F_i)\ne0\) for some \(i\). Similarly choose \(g\in A_1\)
with nonzero linear coefficient. Then \(a'=a+gF_i\) has the same
values as \(a\) on solutions of the equations and at the origin,
and \(\phi(a')=\phi(g)\phi(F_i)\ne0\). If \(\phi(a)\ne0\) already,
take \(a'=a\). This proves the assertion used in choosing
\(\mathcal J\) in the Rees argument.

For a monomial \(X^\nu\) put
\[
 n(\nu)=\sum_{w_i<0}|w_i|\nu_i,\qquad
 p(\nu)=\sum_{w_i>0}w_i\nu_i .
\]
Its weight is \(p(\nu)-n(\nu)\). The filtration given by the
inequality \(n(\nu)\ge k\) is multiplicative, since \(n\) is
additive on monomials. Put \(F_*=F_{r+1}\) in Smith coordinates, and
\[
 L_0=A_{-2},\qquad L_1=A_2,\qquad
 L_2=A_0F_*^2+A_{-2}^2A_2,\qquad
 L(d)=L_0^{d_0}L_1^{d_1}L_2^{d_2}.
\]
The first submodule has negative filtration at least two, the second
at least zero. Since a nonzero scalar multiple of \(F_*\) lies in
\((\ker x_A)^2\), its image in the chart has no constant or linear
term. A monomial of weight \(-1\) with negative weight at most one
would be a single weight-\(-1\) variable. Thus \(\phi(F_*)\) has
negative filtration at least two, and every element of \(L_2\)
has filtration at least four. It follows that every
\(h\in L(d)\) has weight \(2(d_1-d_0-d_2)\), and
\(\phi(h)\) has \(n\ge u\), with
\[
                 u=2d_0+4d_2,\qquad v=2d_1+2d_2 .
\]
On the layer \(n=u\) the weight identity forces \(p=v\).

Let \(S_d\) be the finite set of monomials with these two weights
\((u,v)\). The coefficient map from \(L(d)\) to \(F^{S_d}\)
records that first layer. Its image \(H_d\) is an \(\mathcal O\)-lattice
of rank \(\#S_d\), as we now prove.

First, multiplying by \(b\in A_0\) acts on this layer by \(x_A(b)\).
Every nonconstant monomial in \(\phi(b)\) has \(n=p>0\), since
there are no variables of weight zero. Thus \(H_d\) is finitely
generated over \(\mathcal O\) once \(L(d)\) is finitely generated
over \(A_0\). Every weight component is finitely generated over
\(A_0\). To see this, use fixed homogeneous algebra generators and,
for a fixed weight, choose the finitely many minimal exponent
vectors of that weight in coordinatewise order. Their finiteness
is Dickson's lemma. Every other monomial of that weight is a multiple
of a minimal one by a weight-zero monomial. This proves the assertion
for the components \(A_{-2},A_2\), and products and finite sums of
finitely generated submodules are finitely generated. The claimed
finiteness of \(H_d\) follows.

Next we isolate every individual coordinate of \(F^{S_d}\). For a
weight-\(\pm1\) variable, take the homogeneous element separating
its linear axis. Its lowest possible pair \((n,p)\) is \((1,0)\)
or \((0,1)\). All the monomials on that layer are linear and only
the desired coefficient is nonzero. For a weight-\(\pm2\) variable,
axis separation disposes of other linear terms. On its lowest
layer the remaining monomials are products of two unit-weight
variables of the same sign. We have already constructed elements
whose products isolate each such quadratic monomial on this layer.
Subtract suitable \(F\)-multiples of these products to cancel their
coefficients. Clearing all denominators yields a homogeneous element
of \(A_{\pm2}\) with a nonzero coefficient on the desired axis and
no other term on the lowest layer. The argument also supplies such
an element for each product of two same-sign unit-weight variables
by direct multiplication. When two such isolating elements are
multiplied, terms at a larger negative weight cannot contribute to
the least negative layer of the product. On that layer, the product
therefore has just the chosen product monomial with nonzero coefficient.

For \(\nu\in S_d\) partition its negative variables into
\(u/2\) groups of weight \(-2\) and its positive variables into
\(v/2\) groups of weight \(2\). This is possible because the
number of weight-one factors on each side is even after the
weight-two factors are removed. The preceding construction gives
an element isolating \(X^\nu\) in
\[
 A_{-2}^{\,d_0+2d_2}A_2^{\,d_1+d_2}
    = L_0^{d_0}L_1^{d_1}(A_{-2}^2A_2)^{d_2}
    \ \subseteq L(d).
\]
Hence \(F H_d=F^{S_d}\). This finite torsion-free module over a DVR
is free, with the asserted rank.

Define the homogeneous ideal \(\mathfrak q\) of the Rees algebra
\(\mathcal R=\bigoplus_d L(d)T^d\) by requiring each coefficient in
degree \(d\) to have zero layer at \(n=u\). The filtration bound
and its multiplicativity show it is an ideal. At each \(d\), its
component is exactly the kernel of the coefficient map. It follows
that the \(d\)-component of \((\mathcal R/\mathfrak q)/\varpi\)
is \(H_d/\varpi H_d\). Its \(A_0\)-action factors through the
surjection \(x_A:A_0\to\mathcal O\), so its length is
the \(\mathcal O\)-length. This equals \(\#S_d\), as asserted for
every \(d\in\mathbb N^3\), including the boundary of the cone.
This proves the leading-chart Hilbert function translated from
\(\mathrm{ReesGradedPieceSurjection}\).

\subsection{Removing the torus direction from a mixed point}

We now prove the local-length bound used before squaring the equations.
Let \(B=A/I_0\) where \(I_0\) is their homogeneous ideal. Reserve \(I_0\) for
this ideal and \(\ker x_A\) for the augmentation ideal. Write \(D=B_0\).
The two model hypotheses we use are that \(D\) is a finite
\(\mathcal O\)-module and that at every mixed geometric point in
characteristic zero the tangent space to the equations has dimension
at most one. In the source the latter hypothesis states that any
two point derivations annihilating \(I_0\) are linearly dependent.
Since \(A\) is finitely generated over \(A_0\), \(B\) is finitely
generated over \(D\). Let \(k\) be an algebraic closure of \(F\).

Fix a mixed geometric point. There are \(f\in B_{-2}\),
\(g\in B_2\) with \(fg\) nonzero there. For the model in this
application the generators over degree zero have weights
\(\pm1,\pm2,0\). If some element of negative weight has nonzero
value then a monomial among its terms does also, and at least one
negative-weight generator occurring in that monomial has nonzero
value. If it has weight \(-1\) take its square. This supplies \(f\). Positive weights give \(g\) in the same way. Let \(\mathfrak p\) be the kernel
of the geometric point restricted to \(D\). Define the graded ring
\[
     H=(B\otimes_{\mathcal O}F)[(fg)^{-1}] .
\]
Both \(f\) and \(g\) are units in \(H\) and have their original
degrees. We claim that \(H\) is reduced.

Set \(C=H/(f-1)\). It is of finite type over \(F\). At any geometric
point \(\xi\) of this slice, \(\xi(f)=1\) and \(\xi(g)\ne0\), so
\(\xi\) is mixed. The Euler tangent vector on the original graded
algebra sends a homogeneous \(h\) of weight \(w\) to
\(w\,\xi(h)\). The Leibniz rule follows by adding weights. The
homogeneity of the equations makes it a tangent vector on \(B\).
Its value on \(f\) is \(-2\). A tangent vector on the slice pulls
back to a tangent vector on \(B\) vanishing on \(f\). These
two vectors are linearly dependent by the mixed-point hypothesis.
Evaluation on \(f\) shows the slice vector vanishes on \(B\).
It also vanishes on every inverted element by the quotient rule,
hence on \(C\). This holds at every geometric point of the slice.

For completeness base change to \(k\). At a maximal ideal
\(\mathfrak m\) of \(C\otimes_F k\) the residue field is \(k\) by
the Nullstellensatz. The dual of \(\mathfrak m/\mathfrak m^2\) is
the geometric tangent space just shown to vanish. Nakayama's lemma
in the Noetherian local ring now gives \(\mathfrak m_{\mathfrak m}=0\).
Each such local ring is a field. A nilpotent which survives in the
algebra survives at some maximal ideal containing its annihilator,
which would contradict the field property there. Thus \(C\otimes_F k\)
and, by faithful scalar extension, \(C\) are reduced.

The assertion lifts from the slice to the open graded ring without
any descent assumption on weight spaces. Introduce an invertible
variable \(U\) with \(U^2=f^{-1}\). There is an isomorphism
\[
 H[U,U^{-1}]/(U^2-f^{-1})\ \simeq\ C[V,V^{-1}]
\]
sending homogeneous \(h\in H_w\) to \(\bar h V^w\) and \(U\) to \(V\).
Its inverse sends \(\bar h\) for a homogeneous representative to
\(hU^{-w}\), and sends \(V\) to \(U\). Indeed \(fU^2=1\) makes
the map factor through the slice, and checking on homogeneous
elements and the invertible variable proves the two compositions
are the identities. The left-hand ring is free of rank two over
\(H\) by the monic equation and invertibility of its constant
coefficient. The right-hand ring is reduced since \(C\) is reduced.
The inclusion of \(H\) is injective, giving the claim.

Now localize \(B\) at \(D\setminus\mathfrak p\), obtaining a graded
ring \(B'\). This factors through \(H\), since all nonzero scalars
of \(\mathcal O\) and \(fg\) lie outside \(\mathfrak p\). Thus \(B'\) is a further localization of \(H\).
Thus \(B'\) is reduced. Its degree-zero ring \(D_{\mathfrak p}\)
is local, reduced, and a localization of the finite-dimensional
\(F\)-algebra \(D\otimes_{\mathcal O}F\). It is therefore a field.

For an arbitrary integer \(w\), the element \(h_*=f^w\) is a
unit of degree \(-2w\) in \(B'\). If \(B'_w\ne0\), choose
\(0\ne y\in B'_w\). Reducedness makes \(y^2 h_*\) nonzero in
the field \(D_{\mathfrak p}\). For each \(z\in B'_w\) the identity
\[
              z=\frac{zyh_*}{y^2 h_*}\,y
\]
then expresses \(z\) as a scalar multiple of \(y\). Hence every
weight piece has dimension at most one over \(D_{\mathfrak p}\).
Pullback along the surjection from the localization of \(A_0\)
does not change its module length. This proves the generic
local-length-one bound at every restriction of a mixed geometric
point. It supplies the input to the squaring argument in
\eqref{eq:single-rees}. The corresponding source proofs are
\(\mathrm{SliceHolds}\) and
\(\mathrm{LocLenQuotientWeightPieceLeOne}\).

The arithmetic construction of the graded model, its chart and the
mixed-point tangent bound are inputs to this section. The lower
Euler estimate for extracting scalar factors, the support comparison
for the large Rees pieces and the global automorphy statements in
Section~\ref{sec:single-character} are further steps of the original
argument.
